\scanpage{74}
\mspage{25}
\unsure{Revenons|Revenant} pour un instant à $X$, et considérons \unsure{l'ens} \deleted{\unsure{de}} $I_0$ des $i\in I$ tel que $X_i=\varnothing$. C'est un crible, et on a $X_i^*=\varnothing$, $V_{ij}^*=V_{ij}=\varnothing$ si $i\in I_0$. Donc négligeant les espaces vides, on voit que le diagramme de type $\widetilde I$ défini par l'espace stratifié $X$ ne change pas \unsure{ess\textsuperscript{t}.} en remplaçant $I$ par $I\setminus I_0$, ou plus gén par $I\setminus I'_0$, où $I'_0\subset\unsure{I'|I_0}$ est un crible, ce qui donne lieu à un diagramme $\widetilde{I\setminus I'_0}$ qui est \uline{contenu} dans $\widetilde I$ (cela est vrai pour \uline{tt} crible de $I$).

Si p.ex. on a deux cribles
\[
(\unsure{14})\quad I''\subset I'\subset I
\]
d'où
\[
(\unsure{15})\quad X_{I''}\subset X_{I'}\quad\text{et}\quad U_{I',I''}\overset{\mathrm{def}}{=}X_{I'}\setminus X_{I''},
\]
on peut \deleted{\unsure{sur}} (et on va) sur $U_{I',I''}$ regarder plutôt la \deleted{filtration} \inserted{stratification} de type $I'\setminus I''$, définie par les
\deleted{\[
X_{i'}\setminus X_{i'}\cap X_{I''}\quad(i'\in I')\subset X
\]}
\[
(\unsure{16})\quad X_{i',\deleted{\unsure{X}}\inserted{I',\unsure{I''|U''}}}=X_{i'}\setminus X_{i'}\cap X_{I''}\subset U_{I',I''}
\]
dont les topos élémentaires sont donc les $X_{i'}^*$ $(i'\in I'\setminus I'')$ et des $V^*_{i',j',X'}$, $V_{i',j',X'}$

\mspage{73}
